\begin{theorem}[Chebyshev’s Sum Inequality]
Let  
\[
a_{1}\le a_{2}\le\cdots\le a_{n},\qquad
b_{1}\le b_{2}\le\cdots\le b_{n}
\]
be real numbers.  Then  
\[
n\sum_{i=1}^{n}a_{i}b_{i}\;\ge\;
\Bigl(\sum_{i=1}^{n}a_{i}\Bigr)\Bigl(\sum_{i=1}^{n}b_{i}\Bigr). \tag{1}
\]
\end{theorem}

\begin{proof}
\textbf{Degenerate cases.}  
If \(n=1\) the inequality reduces to \(a_{1}b_{1}\ge a_{1}b_{1}\), true with
equality.  If all the \(a_{i}\) are equal (or all the \(b_{i}\) are equal)
both sides of \((1)\) equal \(a_{1}\sum_{i}b_{i}\) (respectively
\(b_{1}\sum_{i}a_{i}\)), so equality holds.  These situations will be
mentioned again at the end of the proof.

\medskip\noindent\textbf{Reformulation of the inequality.}  
The right–hand side of \((1)\) can be written as a double sum:
\[
\Bigl(\sum_{i=1}^{n}a_{i}\Bigr)\Bigl(\sum_{i=1}^{n}b_{i}\Bigr)
   =\sum_{i=1}^{n}a_{i}b_{i}
    +\sum_{1\le i<j\le n}\bigl(a_{i}b_{j}+a_{j}b_{i}\bigr). \tag{2}
\]
Subtracting \(\sum_{i}a_{i}b_{i}\) from both sides of \((1)\) and using
\((2)\) we obtain the equivalent statement
\[
(n-1)\sum_{i=1}^{n}a_{i}b_{i}
   \;\ge\;
\sum_{1\le i<j\le n}\bigl(a_{i}b_{j}+a_{j}b_{i}\bigr). \tag{3}
\]

\medskip\noindent\textbf{A non‑negative elementary inequality.}  
For any pair \(i<j\) the monotonicity \(a_{i}\le a_{j}\) and \(b_{i}\le b_{j}\)
implies
\[
(a_{j}-a_{i})(b_{j}-b_{i})\ge 0. \tag{4}
\]
Expanding the product gives
\[
a_{j}b_{j}+a_{i}b_{i}-a_{j}b_{i}-a_{i}b_{j}\ge 0,
\]
or equivalently
\[
a_{i}b_{i}+a_{j}b_{j}\;\ge\;a_{i}b_{j}+a_{j}b_{i}. \tag{5}
\]

\medskip\noindent\textbf{Summation over all unordered pairs.}  
Summing \((5)\) over all \(\binom{n}{2}\) pairs \((i,j)\) with \(i<j\) yields
\[
\sum_{1\le i<j\le n}\bigl(a_{i}b_{i}+a_{j}b_{j}\bigr)
   \;\ge\;
\sum_{1\le i<j\le n}\bigl(a_{i}b_{j}+a_{j}b_{i}\bigr). \tag{6}
\]

Each diagonal term \(a_{k}b_{k}\) appears exactly \(n-1\) times on the left
hand side of \((6)\) (once together with every other index), so
\[
\sum_{1\le i<j\le n}\bigl(a_{i}b_{i}+a_{j}b_{j}\bigr)
   =(n-1)\sum_{k=1}^{n}a_{k}b_{k}. \tag{7}
\]

Substituting \((7)\) into \((6)\) gives precisely inequality \((3)\):
\[
(n-1)\sum_{i=1}^{n}a_{i}b_{i}
   \ge
\sum_{1\le i<j\le n}\bigl(a_{i}b_{j}+a_{j}b_{i}\bigr).
\]

Thus \((3)\) holds, and by the equivalence established in Step 2 we obtain
the original Chebyshev inequality \((1)\).

\medskip\noindent\textbf{Equality condition.}  
Equality in \((4)\) occurs exactly when \((a_{j}-a_{i})(b_{j}-b_{i})=0\).
Hence for every pair \(i<j\) we must have either \(a_{i}=a_{j}\) or
\(b_{i}=b_{j}\).  Because the sequences are non‑decreasing, the only way
all such products vanish is that **one entire sequence is constant**:

\begin{itemize}
\item either \(a_{1}=a_{2}= \dots =a_{n}\), or
\item \(b_{1}=b_{2}= \dots =b_{n}\).
\end{itemize}
The case \(n=1\) is automatically covered (the condition is vacuous).
No other configuration can make all pairwise products zero while preserving
the ordering, so these are the complete equality cases.

\medskip\noindent\textbf{Conclusion.}  
We have rewritten Chebyshev’s sum inequality as a sum over unordered pairs,
used the elementary non‑negativity \((a_{j}-a_{i})(b_{j}-b_{i})\ge0\) that
follows from the similarly ordered sequences, and summed over all index
pairs.  The argument works for every \(n\ge1\).  Equality holds **iff** at
least one of the sequences is constant (or trivially when \(n=1\)). \(\square\)
\end{proof}